% Original source: ne630_problems/lesson_04/statement.tex @ e1fba66, lines 92--241.
% Solution prose retained; one documented sign typo is corrected below.
\noindent{\bf Problem 1 [FNRP 2.2]}. The uncollided flux at a distance $r$ from a point source emitting neutrons is given by Eq.~(2.9).

\begin{enumerate}[label=(\alph*)]
\item If you are 1 m away from a very small 1-Ci source of neutrons, what is the flux of neutrons in n/cm$^2$/s, neglecting scattering and absorption in air?
\item If a (1-m thick) shield is placed between you and the source, what absorption cross section would be required to reduce the flux by a factor of 10?
\item Suppose the shield made of the material specified in part (b) is only 0.5 m thick.  How far must you be from the source for the flux to be reduced by the same amount as in (b)?
\end{enumerate}


\vspace{0.5cm}
\ifsolved
{
\color{purple}

\noindent{\bf SOLUTION}:


\begin{enumerate}[label=(\alph*)]
\item
Because the source is ``very small,'' it can be treated as a point source.  Without attenuation due to collisions with air, the neutron flux is given by
\begin{equation*}
 \phi_u = \frac{s_p}{4\pi r^2} = \frac{3.7\times 10^{10}~\text{n}\per\second}{4\pi (100)^2~\centi\meter^2}
 \approx \boxed{2.94 \cdot 10^{5} \frac{\text{n}}{\centi\meter^2~\second}} \, .
\end{equation*}
\item By placing a 1-m thick shield of material with an absorption cross section equal to $\Sigma_a$ between you and the source, the flux is attenuated by a factor equal to $e^{-\Sigma_a\cdot 100}$.  We want this attenuation to be equal to $0.1$, so
\begin{align*}
  e^{-\Sigma_a\cdot 100} &= 0.1 \\
  % Correction: the leading minus sign was omitted in the legacy source.
  \Aboxed{\Sigma_a  &= -\frac{\ln(0.1)}{100} \approx 0.023~\frac{1}{\centi\meter}} \, .
\end{align*}
\item If we have only 50 cm of shielding, the uncollided flux at a distance $r > 50$ from the point source is
\begin{equation*}
 \phi_u(r) = \frac{s_p e^{-\Sigma_a \cdot 50}}{4\pi r^2}  \, .
\end{equation*}
Set this equal to the flux value of (a) times the attenuation of (b), or
\begin{align*}
 \frac{s_p e^{-\Sigma_a \cdot 50}}{4\pi r^2} &= \frac{s_p e^{-\Sigma_a \cdot 100}}{4\pi 100^2} \\
 \frac{1}{r^2} &= \frac{e^{-\Sigma_a \cdot 50}}{100^2} \\
 \frac{1}{r^2} &= \frac{e^{-0.023 \cdot 50}}{100^2} \\
 \Aboxed{r &\approx 178~\centi\meter} \, .
\end{align*}


\end{enumerate}

}
\newpage
\else
% NOTHING
\fi

%%%%%%%%%%%%%%%%%%%%%%%%%%%%%%


\noindent{\bf Problem 2 [FNRP 2.4]}. A boiling water reactor operates at 1000 psi.  At that pressure, the density of water and of steam are, respectively, 0.74 g/cm$^3$ and 0.036 g/cm$^3$.  The microscopic cross sections of H and O are 21.8 b and 3.8 b.
\begin{enumerate}[label=(\alph*)]
\item What is the macroscopic total cross section of the water?
\item What is the macroscopic total cross section of the steam?
\item If, on average, 40\% of the volume is occupied by steam, then what is the macroscopic total cross section of the steam water mixture?
\item What is the macroscopic total cross section of water under atmospheric conditions at room temperature?
\end{enumerate}

\vspace{0.5cm}
\ifsolved
{
\color{purple}

\noindent {\bf SOLUTION}:

\begin{enumerate}[label=(\alph*)]
\item The number density of water molecules is given by
\begin{equation*}
  n_{w} = \frac{\rho_w N_a}{M_w} = \frac{0.74 \cdot 6.022\cdot 10^{23}}{2\cdot 1 + 16}
   \approx 2.4757\cdot 10^{22}~\frac{\text{molecules}}{\centi\meter\cubed}
\end{equation*}
The number density of oxygen is $n_{\text{O}} = n_w$, while $n_{\text{H}} = 2n_w$.  Thus, the macroscopic cross section of water is
\begin{align*}
 \Sigma_w &= n_{\text{O}}\sigma^{\text{O}} + n_{\text{H}}\sigma^{\text{H}}\\
          &= 2.4757\cdot 10^{22} \frac{\text{a}}{\centi\meter\cubed} ~
             ( 3.8~\frac{\text{b}}{\text{a}} + 2\cdot 21.8~\frac{\text{b}}{\text{a}} )
               \cdot 10^{-24}~\frac{\text{\centi\meter\squared}}{\text{b}} \\
          &\approx \boxed{1.17~\frac{1}{\centi\meter}} \, .
\end{align*}


\item Because the number density and, hence, macroscopic cross section scale with the mass density, the macroscopic cross section of steam is
\begin{align*}
 \Sigma_s &= \frac{\rho_s}{\rho_w} \Sigma_w \\
          &= \frac{0.036}{0.74} \cdot 1.17 \\
          &\approx \boxed{0.057~\frac{1}{\centi\meter}} \, .
\end{align*}

\item  Given volume fractions, we can compute an average number density for H and O and recompute a $\Sigma$.  However, that is equivalent to weighting $\Sigma_w$ and $\Sigma_s$, or
\begin{equation*}
 \Sigma_{\text{mix}} = 0.4\Sigma_{s} + 0.6\Sigma_{w} \approx \boxed{0.72~\frac{1}{\centi\meter}} \, .
\end{equation*}
That the cross section depends so strongly on the steam quality makes accurate descriptions of the water coolant in reactors very important.

\item At standard temperature and pressure,
$\rho_w \approx 1~\gram\per\centi\meter\cubed$, so the modified cross section
is
\[
 \boxed{\Sigma_{w'}=(1/0.74)\Sigma_w
 =1.58~\centi\meter^{-1}}.
\]
\end{enumerate}


} %%%%%%%%%%%%%%%%%%
\newpage
\else
% NOTHING
\fi
%%%%%%%%%%%%%%%%%%%%%%%%%%%%%%


\noindent{\bf Problem 3 [FNRP 2.8]}.
What is the total macroscopic thermal cross section of uranium dioxide (UO$_2$) that has been enriched to 4\%?  Assume $\sigma^{235} = 607.5$ b, $\sigma^{238} = 11.8$ b, $\sigma^{O} = 3.8$ b, and that UO$_2$ has a density of 10.5 g/cm$^3$.


\vspace{0.5cm}
\ifsolved
{
\color{purple}

\noindent {\bf SOLUTION}:

First, we assume that the enrichment is by {\it atomic} composition, which is consistent with the book (see the top of page 36).  Then, the molecular mass of 4\% enriched UO$_2$ is
\begin{equation*}
 M_{\text{UO}_2} = 0.04\cdot 235 + 0.96\cdot 238 + 2\cdot 16 \approx 269.88~\frac{\gram~\text{UO}_2}{\text{mol}} \, ,
\end{equation*}
where the mass number of each nuclide is used as the molecular mass (instead of tabulated masses).  This introduces a small but negligible error for these sorts of calculations.  We can simplify further, too, by assuming $M_U = 238$ (i.e., we assume the same mass for both isotopes) for which $M_{\text{UO}_2} = 270$.  I'll keep the first value in my solution.

The number density of UO$_2$ molecules is given by
\begin{align*}
  n_{\text{UO}_2} &= \frac{\rho_{\text{UO}_2} N_a}{M_{\text{UO}_2}} \\
                   &= \frac{10.5 \cdot 6.022\cdot 10^{23}}
                           {269.88} \\
    &\approx 2.3429\cdot 10^{22}~\frac{\text{molecules}}{\centi\meter\cubed}
\end{align*}
Then $n_{235} = 0.04 n_{\text{UO}_2}$, $n_{238} = 0.96 n_{\text{UO}_2}$, and $n_{\text{O}}=2 n_{\text{UO}_2}$, so that the macroscopic cross section is
\begin{align*}
 \Sigma^{\text{UO}_2}_t
  &= n_{235} \sigma_t^{235} + n_{238} \sigma_t^{238} + n_{\text{O}} \sigma_t^{\text{O}} \\
  &= n_{\text{UO}_2} ( 0.04 \sigma_t^{235} + 0.96 \sigma_t^{238} + 2 \sigma_t^{\text{O}}) \\
  &= 2.3429\cdot 10^{22} ( 0.04(607.5) + 0.96(11.8) + 2(3.8) )\cdot 10^{-24} \\
  &\approx \boxed{1.01~\frac{1}{\centi\meter}} \, .
\end{align*}




} %%%%%%%%%%%%%%%%%%
\newpage
\else
% NOTHING
\fi
